\section{Pradinių svorių pasirinkimas}
\label{sec:pradiniu_svoriu_pasirinkimas}

Pradinės svorių (angl.~\textit{weights}) ir poslinkio (angl.~\textit{bias}) reikšmės buvo parinktos atsitiktinai iš tolygaus skirstinio intervale $[-0{,}5;\ 0{,}5]$. Kadangi duomenų aibėje yra $9$ požymiai, tai buvo sugeneruoti $9$ svoriai $w_1, w_2, \dots, w_9$ ir vienas poslinkis $w_0$.

Maži atsitiktiniai pradiniai svoriai yra svarbūs dėl kelių priežasčių:
\begin{itemize}
    \item Jei visi svoriai būtų vienodi, neuronas negalėtų efektyviai mokytis, nes visi gradientai būtų simetriški.
    \item Jei pradinės reikšmės būtų per didelės, sigmoidinės funkcijos~\eqref{eq:sigmoid} išėjimas patektų į sritis arti $0$ arba $1$, kur funkcijos kreivė yra beveik plokščia. Tokiu atveju gradientai tampa labai maži ir mokymas praktiškai sustoja.
    \item Maži atsitiktiniai svoriai užtikrina, kad sigmoidinės funkcijos~\eqref{eq:sigmoid} išėjimas pradžioje yra arti $0{,}5$, kur gradientai yra didžiausi ir mokymas vyksta efektyviausiai.
\end{itemize}

Atsitiktinių skaičių generatoriaus pradinė reikšmė (angl. \textit{seed}) buvo fiksuota ($seed = 42$), kad eksperimentų rezultatai būtų atkartojami.